\chapter{Cor e Espaço de Cor}

A percepção de cor é um dos aspectos mais fundamentais e, ao mesmo tempo, um dos mais complexos da computação gráfica moderna. Longe de ser apenas um trio de intensidades $[R, G, B]$ em um shader, a cor representa o resultado final de uma intrincada cadeia de interações físicas, biológicas e psicológicas. Como afirma \textcite{shirley2021fundamentals}, para renderizar imagens convincentes, o desenvolvedor deve tratar a cor não como uma propriedade absoluta dos objetos, mas como um fenômeno dinâmico que depende da fonte de luz, das propriedades ópticas dos materiais e da sensibilidade do observador humano. Este capítulo detalha como a radiação eletromagnética se torna cor e como os motores gráficos modernos gerenciam esse processo através de pipelines de iluminação linear, HDR e técnicas avançadas de gerenciamento de cor.

\section{A Física da Luz e da Cor}

A luz é uma forma de radiação eletromagnética que apresenta um comportamento dual de onda-partícula, conforme descrito pela mecânica quântica. Do ponto de vista da renderização, estamos interessados principalmente na sua natureza ondulatória e na sua distribuição de energia através do espectro.

\subsection{O Espectro Visível}

O espectro eletromagnético é vasto, abrangendo desde ondas de rádio até raios gama, mas o sistema visual humano é sensível apenas a uma faixa extremamente estreita. Os comprimentos de onda visíveis situam-se aproximadamente entre $380$ nm (limite do violeta) e $780$ nm (limite do vermelho). 

\begin{itemize}
    \item \textbf{Monocromatismo:} Luzes que contêm apenas um comprimento de onda são chamadas de cores espectrais puras. Elas são raras na natureza, ocorrendo principalmente em fenômenos de dispersão como o arco-íris ou através de dispositivos artificiais como lasers.
    \item \textbf{Espectro Contínuo:} A luz solar e as lâmpadas incandescentes emitem uma mistura de quase todos os comprimentos de onda visíveis em intensidades variadas. Essa mistura é descrita matematicamente por uma Distribuição de Potência Espectral (SPD), que mapeia a energia em função do comprimento de onda $\lambda$.
\end{itemize}

A cor não existe no mundo físico de forma independente da percepção. O que existe é a radiação com diferentes frequências. A cor é uma interpretação subjetiva gerada pelo sistema nervoso central em resposta aos impulsos elétricos gerados nos olhos.

\section{Percepção Humana de Cor}

O sistema visual humano é o filtro final através do qual toda imagem digital é interpretada. Compreender suas limitações e sensibilidades é crucial para projetar algoritmos gráficos eficientes e realistas.

\subsection{Bastonetes e Cones}

A retina humana possui dois tipos principais de células fotorreceptoras que operam em diferentes regimes de iluminação:
1. \textbf{Bastonetes:} São responsáveis pela visão escotópica, ativa em condições de baixíssima luminosidade. Eles são extremamente sensíveis ao brilho, permitindo-nos detectar formas e movimentos no escuro, mas são incapazes de distinguir comprimentos de onda, o que resulta em uma visão acromática (tons de cinza).
2. \textbf{Cones:} São responsáveis pela visão fotópica (luz do dia) e pela percepção cromática. Existem três tipos de cones em um olho humano padrão, cada um com pigmentos sensíveis a diferentes partes do espectro:
\begin{itemize}
    \item \textbf{S (Short):} Sensível a comprimentos de onda curtos, com pico de absorção em aproximadamente $420$ nm (Azul).
    \item \textbf{M (Medium):} Sensível a comprimentos de onda médios, com pico em $530$ nm (Verde).
    \item \textbf{L (Long):} Sensível a comprimentos de onda longos, com pico em $560$ nm (Vermelho).
\end{itemize}

\subsection{Teoria Tricromática e o Fenômeno do Metamerismo}

A existência de apenas três tipos de cones implica que a visão humana realiza uma redução drástica de dados. Qualquer SPD complexa, composta por infinitos comprimentos de onda, é projetada em apenas três valores de estimulação neural ($S, M, L$), conhecidos como valores tristimulus.

Este fato biológico dá origem ao \textbf{Metamerismo}: duas fontes de luz com SPDs completamente distintas podem ser percebidas como a mesma cor exata se a estimulação que elas provocam nos cones for idêntica. Este princípio é a pedra angular da tecnologia de displays; ele permite que monitores combinem apenas três primários (Vermelho, Verde e Azul) para simular quase qualquer cor perceptível pelo olho humano, mesmo sem reproduzir o espectro original da luz refletida pelos objetos reais.

\subsection{Sensibilidade Seletiva e Implicações na Engenharia}

A sensibilidade do olho humano não é uniforme através do espectro. Devido à evolução, temos uma densidade muito maior de cones M e L do que de cones S. Além disso, os picos de sensibilidade dos cones M e L estão muito próximos. Isso nos torna significativamente mais sensíveis a variações na região do verde e do amarelo do que na região do azul. 

Essa característica biológica dita padrões fundamentais na engenharia de sensores e compressão:
\begin{enumerate}
    \item \textbf{Bayer Pattern:} Os sensores de câmeras digitais utilizam o dobro de pixels verdes em relação aos vermelhos ou azuis para mimetizar a maior resolução espacial do olho humano nessa faixa cromática.
    \item \textbf{Formatos de Pixel:} Formatos de cor históricos como o R5G6B5 dedicam 6 bits ao canal verde e apenas 5 aos demais, priorizando a precisão onde ela é mais notada.
    \item \textbf{Luminância:} Ao converter uma cor para um valor de brilho, o canal verde recebe o maior peso matemático.
\end{enumerate}

\section{Metrologia de Cor: CIE XYZ e o Diagrama de Cromaticidade}

Para quantificar a cor de forma científica, a \textit{Commission Internationale de l'Éclairage} (CIE) estabeleceu em 1931 o espaço de cor XYZ. Este sistema é baseado em funções de combinação de cores padronizadas $\bar{x}(\lambda), \bar{y}(\lambda), \bar{z}(\lambda)$, que representam a resposta média de um observador humano.

\subsection{Derivação dos Valores XYZ}

Os valores tristimulus $X, Y, Z$ de uma fonte de luz com distribuição de potência espectral $S(\lambda)$ são calculados através de integrais sobre o espectro visível:

\begin{formula}{Cálculo dos Valores Tristimulus CIE XYZ}
\begin{align}
    X &= \int_{380}^{780} S(\lambda)\bar{x}(\lambda)d\lambda \nonumber \\
    Y &= \int_{380}^{780} S(\lambda)\bar{y}(\lambda)d\lambda \nonumber \\
    Z &= \int_{380}^{780} S(\lambda)\bar{z}(\lambda)d\lambda \nonumber
\end{align}
Onde $Y$ é definido propositalmente para corresponder à luminância percebida.
\end{formula}

\subsection{O Diagrama de Cromaticidade}

Para visualizar as cores independentemente do seu brilho, utilizamos as coordenadas de cromaticidade $x, y, z$, que normalizam os valores tristimulus:

\begin{equation}
    x = \frac{X}{X+Y+Z}, \quad y = \frac{Y}{X+Y+Z}, \quad z = \frac{Z}{X+Y+Z}
\end{equation}

Como $x+y+z=1$, basta plotar $x$ e $y$ para representar todas as cores possíveis no famoso "diagrama de ferradura" (horseshoe diagram). A borda curva do diagrama representa as cores espectrais puras, enquanto a linha reta na base representa a "linha das púrpuras". 

\subsubsection{Locus Planckiano e Temperatura de Cor}

Dentro do diagrama de cromaticidade, o \textbf{Locus Planckiano} é a curva que representa as cores emitidas por um corpo negro ideal à medida que sua temperatura aumenta. Esta curva atravessa a região central do diagrama, partindo do vermelho profundo (baixas temperaturas) em direção ao azul (altas temperaturas). A temperatura de cor de uma luz é definida pelo ponto nesta curva que mais se aproxima da cromaticidade da luz em questão.

\section{O Modelo de Cor RGB}

O modelo RGB (Red, Green, Blue) é um sistema de cor aditivo, fundamentado na soma de luzes primárias. Ele pode ser visualizado como um espaço vetorial tridimensional, onde cada cor possível é representada como um ponto dentro de um cubo unitário.

\subsection{Luminância Relativa}

A luminância ($Y$) é uma medida fotométrica da intensidade luminosa por unidade de área, ponderada pela sensibilidade espectral humana. Para o padrão de monitores modernos e sinal de vídeo de alta definição (ITU-R BT.709), a fórmula de luminância relativa é:

\begin{formula}{Luminância Relativa Rec. 709}
\begin{equation}
    Y = 0.2126R + 0.7152G + 0.0722B
\end{equation}
\end{formula}

Observe que o verde contribui com aproximadamente $71,5\%$ da percepção de brilho, enquanto o azul contribui com apenas $7,2\%$. Esta fórmula é essencial para tarefas como a conversão de imagens para escala de cinza ou para o cálculo de efeitos de pós-processamento que dependem do brilho local, como o Bloom.

\section{Modelos de Cor Intuitivos: HSV e HSL}

Enquanto o RGB é ideal para o controle eletrônico de hardware, modelos cilíndricos são muito mais intuitivos para a criação artística e seleção de cores.

\subsection{HSV (Hue, Saturation, Value)}

O modelo HSV organiza as cores de forma que o matiz (Hue) seja um ângulo em um círculo cromático. A saturação representa a "pureza" da cor, e o valor (Value) representa a intensidade máxima de qualquer um dos canais primários.

\begin{align}
    V &= \max(R, G, B) \nonumber \\
    \Delta &= V - \min(R, G, B) \nonumber \\
    S &= \begin{cases} 0, & \text{se } V = 0 \\ \frac{\Delta}{V}, & \text{caso contrário} \end{cases} \nonumber
\end{align}

\subsection{HSL (Hue, Saturation, Lightness)}

O HSL é semelhante, mas define a luminosidade ($L$) como a média entre os componentes RGB máximo e mínimo. Isso coloca o branco puro no topo do cilindro ($L=1$), enquanto no HSV o branco é alcançado com valor máximo e saturação zero.

\section{Gamma e o Fluxo de Trabalho Linear}

Esta é a seção técnica mais crítica para a fidelidade de qualquer motor gráfico moderno. A relação entre a voltagem de entrada de um monitor e a luz efetivamente emitida não é linear, mas segue uma curva de potência.

\subsection{O Princípio de Weber-Fechner e Gamma}

A percepção humana de brilho é logarítmica. Somos muito mais sensíveis a pequenas variações em tons escuros do que em tons claros. Se armazenássemos cores de forma linear em 8 bits por canal (256 níveis), sofreríamos com o efeito de \textbf{banding} nas sombras, enquanto desperdiçaríamos enorme precisão nas altas luzes.

\subsection{A Curva sRGB}

Para otimizar o armazenamento e a transmissão de imagens, aplicamos uma função de transferência chamada \textbf{Gamma Encoding}. O padrão sRGB utiliza uma curva que dedica mais valores numéricos aos tons escuros, alinhando-se à sensibilidade humana.

\begin{formula}{Aproximação da Curva Gamma}
A relação simplificada entre o valor linear ($C_{lin}$) e o valor sRGB armazenado ($C_{srgb}$) é dada por:
\begin{equation}
    C_{srgb} = C_{lin}^{\frac{1}{2.2}}
\end{equation}
\end{formula}

\subsection{Shading Linear: A Regra de Ouro}

O erro clássico em computação gráfica é realizar cálculos de iluminação diretamente em cores sRGB. Como as leis da física da luz são lineares, os cálculos de BRDF e transporte de luz devem ser feitos obrigatoriamente em espaço linear. Se você iluminar um pixel sRGB de valor $0.5$ e dobrar a intensidade da luz, o resultado visual não será o dobro do brilho percebido.

\textbf{O Ciclo Linear Correto:}
1. \textbf{Entrada:} Todas as texturas que representam cor (Albedo, Emissive) devem ser convertidas para linear ao serem lidas: $c^{2.2}$.
2. \textbf{Processamento:} Cálculos de luz e reflexões ocorrem no espaço linear.
3. \textbf{Saída:} Antes da exibição final, o valor linear deve passar pela correção de gamma: $c^{1/2.2}$.

\begin{lstlisting}[language=GLSL]
// Shader com Linear Workflow
void main() {
    // Leitura e Degamma
    vec3 baseColor = pow(texture(uTex, uv).rgb, vec3(2.2));
    
    // Calculos de Shading (Lineares)
    vec3 light = calculateLight(baseColor);
    
    // Gamma Correction final
    fragColor = vec4(pow(light, vec3(1.0/2.2)), 1.0);
}
\end{lstlisting}

\section{High Dynamic Range (HDR) e Tone Mapping}

No mundo real, a diferença de intensidade luminosa entre uma chama de vela e o brilho do sol é de várias ordens de magnitude. O HDR permite que o motor gráfico use valores muito acima de $1.0$, preservando toda a energia da cena.

\subsection{Tone Mapping: Operadores Clássicos}

Tone mapping é a técnica de mapear o vasto intervalo dinâmico do HDR para o intervalo limitado $[0, 1]$ dos monitores SDR comuns. O operador de Reinhard é um dos mais simples:
\begin{equation}
    T(x) = \frac{x}{1+x}
\end{equation}

\subsection{ACES (Academy Color Encoding System)}

O ACES é o padrão moderno utilizado no cinema e em motores AAA. Ele produz uma curva em "S" que preserva o contraste nas sombras e cria transições suaves nas altas luzes. Uma aproximação comum para o pipeline RRT+ODT do ACES é:

\begin{formula}{Aproximação ACES Filmic Tone Mapping}
Para um valor de entrada HDR $x$, o valor mapeado é:
\begin{equation}
    ACES(x) = \frac{x(ax+b)}{x(cx+d)+e}
\end{equation}
Onde as constantes são $a=2.51, b=0.03, c=2.43, d=0.59, e=0.14$.
\end{formula}

\begin{lstlisting}[language=GLSL, caption=Implementação GLSL do ACES]
vec3 ACESFilm(vec3 x) {
    float a = 2.51;
    float b = 0.03;
    float c = 2.43;
    float d = 0.59;
    float e = 0.14;
    return clamp((x*(a*x+b))/(x*(c*x+d)+e), 0.0, 1.0);
}
\end{lstlisting}

\section{Espaços de Cor para Renderização: ACES AP0/AP1}

Para garantir a maior fidelidade cromática, grandes estúdios utilizam o pipeline ACES completo, que divide o processo em etapas claras:
\begin{itemize}
    \item \textbf{Input Transform (IDT):} Converte o formato de entrada para o espaço de trabalho ACES.
    \item \textbf{ACES AP1 (ACEScg):} É o espaço de cores otimizado para renderização e composição, com primários que cobrem quase todo o espectro visível mas evitam cores negativas. Seus primários são baseados em luzes reais, ao contrário dos primários matemáticos do sRGB.
    \item \textbf{Reference Rendering Transform (RRT):} Aplica a estética visual de filme fotográfico.
    \item \textbf{Output Transform (ODT):} Converte o resultado final para o gamut do monitor específico (sRGB, Rec.709, Rec.2020).
\end{itemize}

O uso do ACES AP1 garante que cores intensas e luzes super-brilhantes se comportem de maneira previsível durante o cálculo de iluminação, evitando o "estouro" de canais individuais que causa artefatos de cor em pipelines sRGB simples.

A padronização do ACES permite que diferentes softwares e motores gráficos mantenham a consistência visual durante todo o processo de produção, desde a captura até a exibição final.

\section{Funções de Mistura de Cor (Color Grading)}

O \textit{Color Grading} é o processo final de estilização da imagem. Em vez de aplicar dezenas de filtros matemáticos complexos para cada pixel, motores modernos utilizam \textbf{Look-Up Tables (LUTs) 3D}.

Uma LUT 3D é uma textura cúbica (geralmente $32 \times 32 \times 32$) onde os eixos RGB originais servem como coordenadas para buscar a nova cor desejada.
O pipeline funciona assim:
1. O artista visualiza a cena e gera um "neutral LUT" (uma imagem que representa a identidade, onde $R_{in}=R_{out}$, etc.).
2. O artista aplica as correções de cor desejadas no Photoshop ou DaVinci Resolve sobre essa imagem LUT.
3. O motor gráfico lê a LUT modificada e a utiliza no fragment shader final para transformar cada pixel da cena original.

\subsubsection{Interpolação em LUTs}

Para evitar artefatos de bandeamento em LUTs de baixa resolução, o shader utiliza \textbf{interpolação trilinear}. No entanto, implementações de alta performance podem utilizar a \textbf{interpolação tetraédrica}, que oferece uma precisão superior ao dividir cada cubo da LUT em seis tetraedros, reduzindo a carga computacional e melhorando a fidelidade cromática nas transições de matiz.

\begin{lstlisting}[language=GLSL, caption=Amostragem de LUT 3D]
vec3 applyLUT(sampler3D lut, vec3 color) {
    float size = float(textureSize(lut, 0).x);
    vec3 scale = vec3((size - 1.0) / size);
    vec3 offset = vec3(1.0 / (2.0 * size));
    return texture(lut, color * scale + offset).rgb;
}
\end{lstlisting}

\section{White Balance e Adaptação Cromática}

O \textbf{Balanço de Branco} ajusta a temperatura de cor da imagem para compensar a cor da luz ambiente. O cérebro humano faz isso automaticamente através da constância de cor, mas os motores gráficos precisam de ajustes explícitos.

\subsection{Transformação de Bradford e von Kries}

A adaptação cromática moderna utiliza o modelo de von Kries, que escala os canais de resposta dos cones. A matriz de Bradford é a mais utilizada para converter cores entre diferentes iluminantes (ex: de D65 para D50).

\begin{formula}{Matriz de Adaptação Cromática de Bradford}
A transformação para o espaço de resposta de cones "espectralmente aguçado" é dada por:
\begin{equation}
M_B = \begin{bmatrix} 0.8951 & 0.2664 & -0.1614 \\ -0.7502 & 1.7135 & 0.0367 \\ 0.0389 & -0.0685 & 1.0296 \end{bmatrix}
\end{equation}
\end{formula}

Para aplicar a transformação, primeiro converte-se do espaço XYZ para o espaço de cones usando $M_B$, ajusta-se o ganho baseado na razão entre as respostas dos iluminantes de origem e destino, e então inverte-se a transformação para retornar ao espaço XYZ.

\begin{lstlisting}[language=GLSL, caption=Implementação simples de Temperatura de Cor]
vec3 applyWhiteBalance(vec3 color, float temperature) {
    // Aproximação para ajuste de temperatura (K)
    // Cores quentes diminuem o azul e aumentam o vermelho
    vec3 factor = vec3(1.0 + temperature, 1.0, 1.0 - temperature);
    return color * factor;
}
\end{lstlisting}

\section*{Exercícios}

\begin{enumerate}
    \item Explique por que a sensibilidade do olho humano ao verde influenciou o design dos sensores Bayer.
    \item Calcule a luminância Rec. 709 para um pixel com cor RGB $[0.2, 0.5, 0.3]$.
    \item Demonstre matematicamente por que a soma de cores em espaço sRGB distorce o resultado visual final.
    \item Implemente em pseudocódigo a função de conversão completa de RGB para HSV.
    \item Compare os operadores de Tone Mapping Reinhard e ACES em termos de preservação de contraste.
    \item O que acontece se aplicarmos a correção de gamma sobre uma imagem de dados, como um normal map?
    \item Descreva a estrutura do formato de arquivo .hdr (RGBE).
    \item Explique o fenômeno do metamerismo e por que ele é essencial para a existência de monitores coloridos.
    \item Como o espaço de cor Rec. 2020 se compara ao sRGB em termos de cobertura do espectro visível?
    \item Derive as coordenadas de cromaticidade $(x, y)$ para uma fonte de luz com valores $X=20, Y=50, Z=30$.
    \item Explique a diferença entre os espaços de cor ACES AP0 e AP1 e qual deles é recomendado para renderização.
    \item Como a interpolação tetraédrica melhora a performance e a qualidade do color grading via LUTs?
\end{enumerate}

\textcite{shirley2021fundamentals}
\textcite{akenine2018realtime}
\textcite{pharr2016pbr}
\textcite{reinhard2010high}
\textcite{hable2010filmic}
\textcite{poynton2012digital}
\textcite{wyszecki1982color}
\textcite{stokes1996standard}

